% PDF de ejemplo de liberadoPDF: 4 páginas con una animación del paquete animate.
% Se genera con scripts/generar_ejemplo.sh (pdfLaTeX). Sin compresión de objetos para que
% las pruebas puedan leer /Count, /PageLabels y el widget de animate como texto: PDF 1.4 (sin
% flujos de objetos), fijado antes de \documentclass porque beamer ya escribe datos al cargarse.
\pdfminorversion=4
\pdfobjcompresslevel=0
\documentclass[aspectratio=169]{beamer}
\usepackage[utf8]{inputenc}
\usepackage[T1]{fontenc}
\usepackage{tikz}
\usepackage{animate}
\definecolor{azul}{HTML}{0077B6}
\setbeamertemplate{navigation symbols}{}
\title{liberadoPDF: PDF de ejemplo}
\subtitle{Cuatro páginas para la prueba de humo}
\author{liberadoPDF}
\date{}
\begin{document}

\begin{frame}
  \titlepage
\end{frame}

\begin{frame}{Animación}
  \centering
  % 24 fotogramas a 12 por segundo: una vuelta completa cada 2 s, en bucle y con arranque automático.
  \begin{animateinline}[autoplay,loop,poster=first]{12}
    \multiframe{24}{iAngulo=0+15}{%
      \begin{tikzpicture}
        \useasboundingbox (-3,-3) rectangle (3,3);
        \fill[azul, rotate=\iAngulo] (-1.5,-1.5) rectangle (1.5,1.5);
      \end{tikzpicture}}
  \end{animateinline}
  \par Un cuadrado que gira (24 fotogramas a 12 por segundo).
\end{frame}

\begin{frame}{Pasos}
  \begin{itemize}
    \item Primer paso: esta página tiene la etiqueta 3.
    \pause
    \item Segundo paso: esta también (misma lámina, otro paso).
  \end{itemize}
\end{frame}

\end{document}
